\section{Validación del modelo programado}

La validación del modelo programado se apoya en tres procedimientos
complementarios.

\subsection{Contraste con resultados analíticos}

Para las estaciones del sistema que admiten un tratamiento analítico se
comparó la utilización simulada con la calculada mediante teoría de colas.
Siendo $\lambda$ la tasa de arribos efectiva a la estación, $E[S]$ el tiempo
medio de servicio y $c$ el número de servidores:

\[
\rho = \frac{\lambda\, E[S]}{c}
\]

Para el puesto de triage, con $\lambda = 9{,}91$ pacientes por hora y
$E[S] = 5{,}33$ minutos, resulta $\rho = 0{,}88$. Para los consultorios, con
$E[S] = 20{,}3$ minutos ponderado por la distribución de colores y $c = 4$,
resulta $\rho = 0{,}84$.

La comparación de estos valores con los de la tabla~\ref{tab:piloto}
constituye la prueba de validación más exigente aplicada al modelo, por
cuanto un desvío significativo señala una inconsistencia en los parámetros o
en la estructura.

\subsection{Validación aparente}

El recorrido de los pacientes sobre el plano fue revisado visualmente durante
las corridas, verificando que la secuencia de desplazamientos se corresponda
con el circuito real: que la espera se produzca en la sala y no en los
consultorios, que los pacientes de ambulancia accedan directamente al Shock
Room y que los consultorios se ocupen de manera equilibrada.

\subsection{Validación con el sistema real}

\begin{pendiente}
Esta subsección debe completarse contrastando las salidas del modelo con el
conocimiento de la profesional del hospital consultada: ¿resultan
verosímiles los tiempos de espera obtenidos? ¿se corresponde la proporción de
internaciones con lo observado en la práctica? Se recomienda documentar esa
conversación como evidencia de validación.
\end{pendiente}

